\chapter{Marco Conceptual}

\section{Desperdicio digital informativo}

El concepto de desperdicio digital se refiere al consumo innecesario de recursos tecnol\'ogicos ---datos, energ\'ia, tiempo de procesamiento y atenci\'on humana--- que no agregan valor al usuario final. En el contexto informativo, el desperdicio ocurre cuando:

\begin{itemize}[nosep]
    \item Un usuario carga m\'ultiples p\'aginas web para obtener informaci\'on que podr\'ia presentarse en un solo formato compacto.
    \item El mismo contenido es procesado y almacenado varias veces por sistemas automatizados.
    \item Los modelos de IA reciben y procesan contenido redundante o de baja calidad.
    \item Las redes sociales exponen a los usuarios a informaci\'on fragmentada, repetida o sin contexto, incentivando el \textit{scroll} continuo.
\end{itemize}

Este desperdicio tiene un costo real: consumo energ\'etico en servidores y dispositivos, emisiones de carbono asociadas a la transmisi\'on y procesamiento de datos, y agotamiento de la capacidad de atenci\'on de las personas.

\section{Green Tech y sostenibilidad digital}

\textbf{Green Tech} (Tecnolog\'ia Verde) es un enfoque que busca dise\~nar, desarrollar y utilizar tecnolog\'ia minimizando su impacto ambiental y promoviendo la sostenibilidad. En el \'ambito del software y los datos, esto implica:

\begin{itemize}[nosep]
    \item Eficiencia en el uso de recursos computacionales.
    \item Reducci\'on de datos innecesarios almacenados y transmitidos.
    \item Optimizaci\'on de consultas a modelos de IA.
    \item Dise\~no de sistemas que promuevan h\'abitos de consumo digital responsable.
\end{itemize}

EcoBrief Bolivia se enmarca en esta corriente al atacar directamente el desperdicio informativo: reduce la cantidad de p\'aginas que un usuario necesita cargar, elimina la redundancia en el procesamiento de contenido y optimiza el uso de modelos de lenguaje.

\section{Inteligencia artificial responsable}

La IA responsable busca que los sistemas de inteligencia artificial sean \'eticos, transparentes, eficientes y sostenibles. En EcoBrief, la IA se utiliza solo despu\'es de que el contenido ha sido filtrado, deduplicado y priorizado. Esto significa que:

\begin{itemize}[nosep]
    \item No se env\'ia contenido irrelevante o repetido al modelo.
    \item Se utilizan modelos peque\~nos y r\'apidos para tareas simples.
    \item Los resultados se cachean para evitar reprocesamiento.
    \item Se limita el n\'umero de candidatos por categor\'ia.
\end{itemize}

De esta forma, la IA se convierte en una herramienta de reducci\'on, no de multiplicaci\'on de contenido.
